\section{Introduzione al corso}
La prima parte che trattiamo è quella più ingegneristica, un primo svicolamento che 
si adopera dai metodi computazionali classici che sono rigorosi.
Laddove abbiamo gli strumenti adonei per ottenere ottimi risultati usiamo questi.
La natura dei problemi può essere non lineare per cui il calcolo rigoroso è troppo 
dispendioso o impossibile da applicare. Si tentano strade alternative (quando 
la risposta rigorosa è nota per aver un criterio di valutazione).\\

La fonte primaria ispiratrice della computational intelligence. Nasce come proposta
da un iraniano Zadeh, che propone la logica fuzzy, in cui diceva ottimi risultati potevano
essere ottenuti da singole applicazioni.

Il corso si concentra sulla parte più legata a come adoperare in alcuni contesti un 
tipo di ia piuttosto che un altro.

Tecniche di pre processing dei dati.
La flessibilità di ia e come adoperarla in contesti diversi.
Tutte le varie tecniche vengono adoperate e apprezzate nei vari contesti.

Faremo delle simulazioni con Python. Chiameremo l'ia e gli chiederemo di fare cose, 
in modo che ci restituisce un risultato.

Vedremo come trasformare un segnale musicale in un'immagine.

Fuzzy logic, logica come alternativa alla logica booleana, che è basata su true e false.

L'applicazione che faremo prevalentemente è sul segnale musicale, posso usare protocolli 
informatici (MIDI Musical Instrument Digital Interface).

Posso dire di spostare il calcolo da uno analitico chiuso a uno soft
infatti si chiamava soft computing.
Con il passare del tempo tutte le tecniche di natura biologica e naturale sono state
inglobate in questo paradigma, che è diventato più ampio.

Quali problemi trattano questi algoritmi evolutivi? Problemi di ottimizzazione o 
problemi inversi.
Il problema di ottimizzazione significa che è gestito da equazioni governate da alcuni parametro
vorrei sapere la configurazione ottima di quei parametri per ottenere una configurazione ottima.

I problemi inversi sono quelli in cui abbiamo un certo numero di dati e vorremmo 
sapere quale è la configurazione che ha generato quei dati, o che è più vicina 
a quei dati.
Ho un modello matematico che mi da per esempio le previsioni del tempo, questo modello
non è stato attualizzato ai dati che servono a me. All'interno a quelle equazioni
ci sono dei parametri che sono stati messi in modo generico ma che dovrebbero 
essere attualizzati.
Per esempio il modello di una parabola:
$$y=ax^2+bx+c$$
Con a, b, c generici ho il modello, poi piego il modello al caso che mi interessa, 
ovvero cambio i valori di a, b, c in modo da far sì che la parabola approssimi al meglio i dati osservati.

Altri metodi (es. simulated annealing) sono più adatti a problemi di ottimizzazione (esempio lo usano 
molto i fisici), altri (es. reti neurali) sono più adatti a problemi inversi, ma non è detto che non 
possano essere usati anche per l'altro tipo di problema.

Tutti questi algoritmi sono stocastici, c'è una componente stocastica in modo da aumentare la 
capacità di qeuesti algoritmi di performare bene.
Se uso una determinata funzione per trovare per esempio il minimo di una funzione, 
potrei essere bloccato in un minimo locale, se invece uso una componente stocastica, 
posso uscire da questo minimo locale e trovare un minimo globale.

Questi algoritmi vanno visti in funzione delle loro capacità, quindi vanno testati.
Quale algoritmo va bene per piccoli spazi e quale per grandi spazi.

Le più grandi catetegorie con cui vengono classificate sono:
grado di esplorazione (exploration) e grado di sfruttamento (exploitation).

La rete neurale è difatto un problema di ottimizzazione.
La rete neurale dovrebbe emulare il cervello di un vertebrato, ci sono dei dati che 
vengono presi dall'esterno (banca dati).
Per esempio ho in input $y=f(x)$, poi prendo questo dati che passa per il neurone di ingresso
faccio una sfilza di nodi (chiamati neuroni nascosti), l'informazione di input viene
veicolata dal neurone di ingresso a tutti i neuroni nascosti, viene veicolata in modo tale
da avere dei pesi diversi, dando luogo a un uscita, queste uscite vengono dati in pasto
ad altri neuroni di un altro strato nascosto (perché non è nè input nè output).
Ciascuno di questi passaggi ha un altro peso.
Tutta questa informazione viene poi convogliata verso il neurone d'uscita ($y$). 
Quindi tutta la rete è $f(x)$ e in ingresso ho $x$.

Allora provo a mettere i valori sperimentali, all'inizio i pesi sono a caso. Una volta che 
ho tutti i risulati a caso e so quelli che dovrebbero uscire posso applicare un algoritmo
di ottimizzazione per minimizzare l'errore. Una volta applicato l'algoritmo riinserisco 
i valori e itero finché non ottengo i risultati aspettati.

L'insieme dei pesi dei valori sinapsici è lo spazio delle soluzioni.

Reti neurali con memoria (recurrente neural network).
Reti neurali generative.
